% ------------------------------------------------------------------------
% -*-TeX-*- -*-Hard-*- Smart Wrapping
% ------------------------------------------------------------------------


\chapter{評価実験}
\label{cha:4}

本章では，第\ref{cha:3}章の提案手法を検証するための実装，ベンチマークと評価プロトコル，比較条件，統計設計，およびハイパーパラメータを述べる．本章の記述はすべて公開実装（\texttt{scripts/train\_lora\_persona.py}，\texttt{scripts/run\_evolution.py}，\texttt{scripts/run\_eval.py}，\texttt{scripts/run\_eval\_battery.py}，\texttt{scripts/cloud/}）と一致する．未実施の実験結果には言及しない．

\section{実装}
\label{sec:4.1}

\subsection{ベースモデルとQLoRA学習}
\label{sec:4.1.1}

ベースモデルには\textbf{Qwen3-4B-Instruct-2507}\cite{qwen2025qwen3}（Apache-2.0，Hugging Face上でゲートなし公開）を用いる．4B級を選ぶ理由は，（i）同規模帯では「素のdebateが性能を悪化させうる」ことが報告されており\cite{zhang2025debate,cost2026consensus}，チームレベル選択圧がdebateを機能させるかという研究課題の検証に適した困難条件であること，（ii）公表ベンチマーク値（MMLU-Pro 69.6，SuperGPQA 42.8）に改善余地が残ること，（iii）進化ループで多数のLoRAを単一GPU上に同時サービングできる計算規模であることによる．

世代0の3個体（critic / pragmatist / explorer）は，役割別SFTデータ（各60会話例，JSONL形式，チャットテンプレート適用）を用いた\textbf{QLoRA}\cite{dettmers2023qlora}で学習する．設定は以下のとおりである．

\begin{itemize}
\item \textbf{量子化}: ベースモデルを4bit NF4（double quantization有効）でロードし，計算dtypeはGPUに応じてbfloat16（A100等）またはfloat16（T4，compute capability 7.5のためbf16非対応）へ自動フォールバックする．
\item \textbf{LoRA構成}: rank $r = 32$，$\alpha_{\mathrm{LoRA}} = 2r = 64$，dropout 0.05，対象モジュールは全層の\texttt{q\_proj, k\_proj, v\_proj, o\_proj, gate\_proj, up\_proj, down\_proj}の7種．
\item \textbf{最適化}: 学習率$2 \times 10^{-4}$，エポック数3（Vertex AIジョブ投入設定．学習スクリプト単体の既定値は1），デバイスあたりバッチサイズ1，勾配累積16，\texttt{paged\_adamw\_32bit}，gradient checkpointing有効（学習時はKVキャッシュを無効化して競合を回避），packing無効．
\item \textbf{再現性}: グローバルシード1234（Python / NumPy / PyTorch / CUDAを統一固定）．学習コマンド・シード・ハードウェア構成は\texttt{metadata.json}として成果物に同梱する．
\end{itemize}

各役割のペルソナは，SFTにより重み空間に焼き込まれると同時に，推論時にも短い役割記述（例: 批判的検証者「厳密な検証を重視する批判的思考家．反証・例外・境界条件に敏感」）をシステムプロンプトとして与える．

\subsection{vLLM Multi-LoRAサービング}
\label{sec:4.1.2}

全推論（進化ループの連合評価・最終評価とも）はvLLM\cite{kwon2023efficient}のOpenAI互換APIを通じて行う．サーバは\texttt{--enable-lora}，\texttt{--max\_loras 8}，\texttt{--max\_lora\_rank 64}，\texttt{--max-model-len 32768}，tensor並列1で起動し，ベースモデルを1つだけGPUに常駐させたまま，リクエストの\texttt{model}フィールドでLoRAアダプタを問い合わせ単位に切り替える．進化ループで世代ごとに生成される子アダプタは，動的LoRAロードAPI（\texttt{/v1/load\_lora\_adapter}，\texttt{VLLM\_ALLOW\_RUNTIME\_LORA\_UPDATING=True}）でサーバ再起動なしに登録する．

生成パラメータは，進化ループ内ではtemperature 0.7・top-p 0.9・最大512トークン，最終評価ではQwen3の公式推奨に従いtemperature 0.7・top-p 0.8・top-k 20・最大8{,}192トークンとし，いずれもシード指定付きとする（エージェント$\times$ラウンド，およびSelf-Consistencyのサンプル添字ごとに決定的にオフセット．第\ref{cha:3}章\ref{sec:3.2}節）．進化ループと最終評価で生成条件が異なる点は本研究の限界として第\ref{cha:5}章で述べる．クライアントは問題単位で16並列のリクエストを発行し，vLLMサーバ側の連続バッチングでGPU利用率を確保する．API呼び出しは最大3回の線形バックオフ付きリトライで保護する．

\subsection{Vertex AI Spot構成とプリエンプション耐性}
\label{sec:4.1.3}

計算はGCP Vertex AI Custom Trainingの\textbf{Spot（プリエンプティブル）ジョブ}として実行する（プロジェクトのクォータ上，オンデマンドGPU VMは利用不可）．工程ごとの構成を表\ref{tab:4.vertex}に示す．

\begin{table}[htbp]
\caption{Vertex AI Spotジョブの工程別マシン構成}
\label{tab:4.vertex}
\centering
\small
\setlength{\tabcolsep}{4pt}
\begin{tabular}{|p{3.2cm}|p{4.4cm}|p{5.8cm}|}
\hline
工程 & マシン構成 & 理由 \\
\hline
QLoRA学習（世代0） & \texttt{n1-standard-8} + NVIDIA T4 $\times$1（Spot） & 4bit量子化により6GB級で学習可能．fp16フォールバックで対応 \\
\hline
進化ループ／最終評価 & \texttt{a2-highgpu-1g}（A100 40GB $\times$1，Spot，us-central1） & vLLMのMulti-LoRAカーネル（Punica/Triton）がcompute capability $\geq 8.0$を要求しT4では動作しないため \\
\hline
\end{tabular}
\end{table}

成果物（アダプタ・評価JSON・世代ログ）はGCSバケットにCloud Storage FUSEマウント（\texttt{/gcs/...}）経由で保存する．Spotはいつでもプリエンプトされうるため，すべての評価ランナーは\textbf{問題単位の逐次JSONLキャッシュ}を持ち，再実行時には評価済み問題IDをスキップして途中から再開する．最終評価は複数条件を1ジョブにまとめるバッテリードライバで実行し，完了済みエントリのJSONが存在すれば条件ごとスキップする．進化ループも世代ごとに全ログをJSONへ書き出し，任意の世代から再開可能である．

\section{ベンチマークと評価プロトコル}
\label{sec:4.2}

\subsection{ベンチマーク選定}
\label{sec:4.2.1}

最終評価は表\ref{tab:4.bench}に示す3ベンチマークで行う．いずれもHugging Face上でゲートなしに取得できる（認証トークン不要の方針を維持する）．

\begin{table}[htbp]
\caption{最終評価に用いるベンチマーク}
\label{tab:4.bench}
\centering
\footnotesize
\setlength{\tabcolsep}{3pt}
\begin{tabular}{|p{3.1cm}|p{3.1cm}|p{1.2cm}|p{5.6cm}|}
\hline
ベンチマーク & 抽出方法 & 問題数 & 選定根拠 \\
\hline
MMLU-Pro\cite{wang2024mmlupro} & test分割から固定ランダム抽出 & 500 & ベースの公表値69.6で飽和しておらず，MADによる改善報告と接続可能．10択で当て推量の影響が小さい \\
\hline
MATH-500\cite{lightman2023lets,hendrycks2021measuring} Level 4--5 & Level 4以上に限定して抽出 & 200 & 数学推論はMADの改善効果が最大と報告される領域．高難度に限定して天井効果を回避 \\
\hline
SuperGPQA\cite{du2025supergpqa} & 固定ランダム抽出 & 300 & ベースの公表値42.8と最も余地が大きい．大学院水準の広分野知識推論 \\
\hline
\end{tabular}
\end{table}

GSM8K\cite{cobbe2021training}はベースモデルで80--92\%と飽和しているため配管確認（スモークテスト）専用とし，主張には用いない．GPQA-Diamond\cite{rein2023gpqa}はデータセットがHFゲート付きであるため不採用とする．

\subsection{生成型評価と回答抽出}
\label{sec:4.2.2}

評価はlm-evaluation-harness型の選択肢対数尤度方式ではなく，\textbf{生成型（generative）評価}で統一する．solo条件とdebate条件を同一プロトコル・同一抽出ロジックで比較するためには，両者とも自由生成テキストから回答を取り出す必要があるからである．全エージェントに厳密な書式指示
\begin{verbatim}
Think step by step. Then give your final answer on the last line
in exactly this format:
ANSWER: <letter | number | final simplified answer>
\end{verbatim}
を与え，出力末尾の\texttt{ANSWER:}行を正規表現で最優先抽出する．\texttt{ANSWER:}行が欠落した場合に限り，文脈条件付きのフォールバック（選択肢タスクでは``answer is (X)''等のパターン，数値タスクでは最後の数値，数式タスクでは最後の\verb|\boxed{}|）を適用する．数式回答はLaTeX表記の正規化（\verb|\frac|$\rightarrow$\texttt{/}，\verb|\sqrt|$\rightarrow$\texttt{sqrt()}，数値の同値判定 $12.0 = 12$等）を経て比較する．抽出不能な出力は無効票（多数決から除外）とし，単独評価では不正解として扱う．\textbf{すべての比較条件で同一の抽出・正規化・採点コードを共有}し，条件間の差が抽出ロジックの差に混入しないことを保証する．

\subsection{進化ループ内の適応度セット}
\label{sec:4.2.3}

進化ループの選択圧（第\ref{cha:3}章の$v(S)$計測）には，MMLU-Pro test分割（約12,000問）から専用シード（777）で抽出した\textbf{固定100問}を全世代共通で用いる．100問という規模は，tinyBenchmarks\cite{polo2024tinybenchmarks}が示した「キュレートされた100例でフルベンチマークのスコアを平均誤差約2\%で推定できる」という結果に基づく．適応度セットは世代間の\textbf{順位付け専用}であり，その絶対値は主張に用いない．最終評価とはシード系列を分離して抽出するため（最終評価はシード1--3），標本の系統的な共有はないが，母集合が共通のため偶発的な重複は起こりうる（期待重複は500問抽出あたり約4問，1\%未満）．適応度セットへの選択圧の過適合（進化版の「テストへの学習」）が疑われる場合に備え，最終評価がこの100問を含まない形でも集計できるよう問題IDを全ログに記録する．

\section{比較条件}
\label{sec:4.3}

最終評価は表\ref{tab:4.cond}に示す8条件で行う（研究設計書\S 6.2の条件1〜6に，プロンプトペルソナの寄与を分離する条件3$'$と実験2の条件7$^\ddagger$を追加．設計書の条件7＝solo適応度進化A1は未実施）．全条件で同一の問題集合・同一シード・同一抽出ロジックを用い，問題単位で対応づけ（paired）が可能な設計とする．debate条件のラウンド数は$R = 1$に統一する．

\begin{table}[htbp]
\caption{最終評価の比較条件}
\label{tab:4.cond}
\centering
\footnotesize
\setlength{\tabcolsep}{3pt}
\begin{tabular}{|p{0.5cm}|p{2.9cm}|p{3.6cm}|p{6.0cm}|}
\hline
\# & 条件 & 構成 & 目的 \\
\hline
1 & ベースsolo CoT & ベースモデル単体，ペルソナなし & 絶対的な基準線 \\
\hline
2 & \textbf{Self-Consistency@9} & ベースモデルから9サンプル$\rightarrow$多数決\cite{wang2023selfconsistency} & \textbf{生成回数を揃えた非協調ベースライン}．3エージェント$\times$最大3生成（round 0＋2ラウンド）＝9生成に相当．既定$R=1$ではdebate側は6生成/問であり，SC@9はベースライン側に生成回数で有利な保守的設定（実測トークンでもSC@9が総量で上回る; \S\ref{sec:4.6.6}）．「MADは計算量を揃えたSCに勝てない」という批判\cite{smit2024should}への直接の対処 \\
\hline
3 & ベース$\times$3温度サンプリングdebate & 同一ベースモデル3体（ペルソナプロンプトなし，シードオフセットのみ） & サンプリング多様性のみのdebate（Self-MoA型対照） \\
\hline
3$'$ & ベース$\times$3プロンプトペルソナdebate & 同一ベースモデル3体＋役割プロンプト & プロンプトによるペルソナ付与の寄与を分離（重み焼き込みvsプロンプトの対照） \\
\hline
4 & gen-0 LoRAチームdebate & SFT直後の3ペルソナLoRA & 進化前のチーム性能．条件5との差が進化の寄与 \\
\hline
5 & \textbf{進化後チームdebate（主要条件）} & 提案手法（Shapley＋sharing，$G=6$世代）の最終代表チーム & 本研究の主張の中核 \\
\hline
6 & 進化後LoRA solo $\times$3 & 条件5の各個体を単独評価 & チーム効果（協調による上乗せ）と個体改善の分離 \\
\hline
7$^\ddagger$ & 再学習チームdebate（実験2） & 能力保持リプレイで再学習した3ペルソナLoRA＋条件付き更新＋匿名化＋logprob重み付き投票 & 実験1の診断（\S\ref{sec:4.6.2}）に基づく処方の検証．ゲート方式（\S\ref{sec:4.6.3}）で各構成要素を検収した上での最終形 \\
\hline
\end{tabular}
\end{table}

なお研究設計段階ではsolo適応度への置換（A1），fitness sharing無効化（A2），naive交叉（A3）のアブレーションを計画し実行フラグとして実装したが（第\ref{cha:3}章参照），実験1で進化の寄与自体が限定的（\S\ref{sec:4.6.1}）と判明したため，計算資源を実験2（処方の検証）に振り向け，これらは未実施である．この判断の含意は第\ref{cha:5}章で論じる．

\section{統計設計}
\label{sec:4.4}

統計処理はMiller\cite{miller2024adding}の評価統計5原則（標準誤差の併記・クラスタ構造の考慮・リサンプルの扱い・paired分析・検出力設計）に準拠する．

\begin{itemize}
\item \textbf{シードと反復}: 主要比較は$K = 3$シード（1, 2, 3）で反復する．各（シード, 問題）対は両条件で同一の生成条件を共有するため，対応ありデータとして扱う．
\item \textbf{主検定（クラスタ構造の考慮）}: 同一問題への複数シードの回答は相関する反復測定であり，（シード, 問題）対を独立観測とみなしてプールする検定は独立性の仮定を満たさない\cite{miller2024adding}．そこで主解析は\textbf{問題IDをクラスタ単位}とし，問題ごとにシード平均を取った精度差$\Delta_i$に対する\textbf{クラスタ符号反転置換検定}（20{,}000回）でp値を，\textbf{クラスタブートストラップ}（問題単位の再抽出，10{,}000回）で95\%信頼区間を求める．参考値として（シード, 問題）対プールのMcNemar正確検定\cite{dror2018hitchhiker}も併記するが，その p値は反保守的であるため結論には用いない\footnote{初稿では（シード, 問題）プールのMcNemar検定を主解析としていたが，反復測定の相関を無視する点が不適切であるため，問題IDクラスタリングを考慮した解析へ差し替えた．有意性の判定はいずれの解析でも変わらなかった．なお両者は推定対象が僅かに異なる（プールは観測単位，クラスタ解析は問題単位の平均差）ため点推定値も僅かに異なる．}．
\item \textbf{同等性の主張}: 差の検定で有意でないことは同等性の証拠にならないため，同等性を主張する場合は等価マージンを$\pm 2$pt（本研究で意味があるとみなす最小効果量より小さい値）と定義したTOST型判定——クラスタブートストラップの90\%信頼区間が$\pm 2$ptに収まるか——を用いる．
\item \textbf{多重比較補正}: 主要な仮説検定は各実験でプール比較6件とし，\textbf{Holm--Bonferroni法}\cite{holm1979simple}でfamily-wise error rateを5\%に制御する．ベンチマーク別の内訳検定およびそれ以外の比較は探索的分析と位置づける．
\item \textbf{頑健性確認}: 実験2の主検定（条件7$^\ddagger$ vs SC@9）が3シードで境界的（$p=0.053$）となったため，\textbf{効果の維持を確認する目的を事前に明記した上で}seed 4--6を追加し，3シード単独と6シード合算の両方を報告する．このシード追加は結果を見た後の逐次的な標本拡大（optional stopping）に該当するため，6シード合算のp値を確認的検定として扱うには本来逐次検定補正を要する．本研究では補正を行わない代わりにこの経緯を明示し，かつ最終的な結論（条件7$^\ddagger$の有意な劣位）が標本拡大により\textbf{棄却側に不利な方向へ}確定した——すなわちoptional stoppingが偽陽性を利する典型的な構図とは逆である——ことを付記する．
\item \textbf{開発と評価の分離}: 実験2の全ての設計選択（リプレイ比率・集約方式・議論プロトコル）は，最終評価と重複しない検収セット（seed 9）上のゲート（\S\ref{sec:4.6.3}）で確定してから最終評価に進んでおり，最終評価セットへの適合を構造的に排除している．
\item \textbf{検出可能効果量}: 対応あり検定（$\alpha = 0.05$，検出力80\%）の目安は$+10$pt $\rightarrow N \approx 180$，$+7$pt $\rightarrow N \approx 360$，$+5$pt $\rightarrow N \approx 700$である．3ベンチマーク統合（$N = 500 + 200 + 300 = 1{,}000$問，$K=3$シードによる分散低減を含む）では\textbf{$+4$〜$5$pt程度の効果が検出可能}であり，Duらが報告するMADの効果量（7〜15pt）に対して十分な感度を持つ．
\item \textbf{実行記録と再現性の限界}: 全実行についてシード・モデルrevision・vLLMバージョン・生成パラメータ・問題IDリストを結果JSONに記録する．ただし同一シードを指定してもvLLMの生成は並列実行状態に依存して完全には再現しない（\S\ref{sec:4.6.4}）ため，これらの記録が担保するのは監査可能性であり，ビット単位の再現ではない．
\end{itemize}

議論品質のLLM-as-judge評価と世代を通じた行動多様性の推移分析は副次評価として計画したが，測定環境の統制（\S\ref{sec:4.6.4}）に伴う再測定を優先したため未実施である．

\section{ハイパーパラメータ}
\label{sec:4.5}

以下の値はすべて実装のデフォルト値（コマンドライン引数の既定値，およびVertexジョブ投入スクリプトの設定値）と一致する．QLoRA学習の設定を表\ref{tab:4.hp-qlora}に，議論・生成の設定を表\ref{tab:4.hp-debate}に，進化ループの設定を表\ref{tab:4.hp-evolve}に，最終評価の設定を表\ref{tab:4.hp-eval}に，サービング・インフラの設定を表\ref{tab:4.hp-infra}に示す．

\begin{table}[htbp]
\caption{QLoRA学習（世代0）のハイパーパラメータ}
\label{tab:4.hp-qlora}
\centering
\small
\begin{tabular}{|p{5.4cm}|p{7.6cm}|}
\hline
項目 & 値 \\
\hline
ベースモデル & Qwen/Qwen3-4B-Instruct-2507 \\
\hline
量子化 & 4bit NF4，double quantization，計算dtype bf16（T4ではfp16） \\
\hline
LoRA rank $r$ / $\alpha_{\mathrm{LoRA}}$ / dropout & 32 / 64 / 0.05 \\
\hline
対象モジュール & q, k, v, o, gate, up, downの各射影（全層） \\
\hline
学習率／エポック & $2 \times 10^{-4}$ ／ 3（Vertex投入設定; スクリプト既定1） \\
\hline
バッチサイズ／勾配累積 & 1 ／ 16 \\
\hline
最適化器 & \texttt{paged\_adamw\_32bit}（gradient checkpointing有効，packing無効） \\
\hline
SFTデータ & 役割別60会話例$\times$3役割 \\
\hline
シード & 1234 \\
\hline
\end{tabular}
\end{table}

\begin{table}[htbp]
\caption{議論・生成のハイパーパラメータ}
\label{tab:4.hp-debate}
\centering
\small
\begin{tabular}{|p{5.4cm}|p{7.6cm}|}
\hline
項目 & 値 \\
\hline
temperature / top-p / 最大トークン & 0.7 / 0.9 / 512 \\
\hline
議論ラウンド数$R$ & 1（round 0の独立回答＋1更新ラウンド） \\
\hline
集約 & 最終ラウンド多数決（同数はシード付き乱数，無効票は除外） \\
\hline
シードオフセット & エージェント$\times$ラウンド: $\mathrm{seed} \cdot 10^4 + 100i + \rho$；SCサンプル: $\mathrm{seed} \cdot 10^3 + k$ \\
\hline
クライアント並列数／リトライ & 16 ／ 3回（線形バックオフ） \\
\hline
\end{tabular}
\end{table}

\begin{table}[htbp]
\caption{進化ループのハイパーパラメータ}
\label{tab:4.hp-evolve}
\centering
\small
\begin{tabular}{|p{5.4cm}|p{7.6cm}|}
\hline
項目 & 値 \\
\hline
世代数$G$ & 6 \\
\hline
サブ集団構成 & 3役割$\times K = 2$（エリート＋子）$= 6$ LoRA \\
\hline
適応度 & 厳密Shapley値$\times$fitness sharing（\texttt{--fitness-mode shapley}） \\
\hline
適応度セット & MMLU-Pro 100問，シード777（最終評価とシード分離） \\
\hline
sharing $\sigma_{\mathrm{share}}$ & 0.3 \\
\hline
交叉 & $\Delta W$空間ブレンド＋ランダム化SVD再分解（\texttt{delta}） \\
\hline
混合比$\alpha$ & $\mathcal{U}(0.3, 0.7)$から世代・役割ごとに抽出 \\
\hline
SVD & oversampling $+8$，power iteration 4回，rank 32に切詰め \\
\hline
突然変異率$\rho_{\mathrm{mut}}$／幅$\sigma_{\mathrm{mut}}$ & 0.3 ／ 0.02（テンソル標準偏差に対する相対値） \\
\hline
世代0変異体 & $\rho_{\mathrm{mut}} = 1.0$，$\sigma_{\mathrm{mut}} = 0.02$ \\
\hline
進化シード & 1234 \\
\hline
\end{tabular}
\end{table}

\begin{table}[htbp]
\caption{最終評価の設定}
\label{tab:4.hp-eval}
\centering
\small
\begin{tabular}{|p{5.4cm}|p{7.6cm}|}
\hline
項目 & 値 \\
\hline
ベンチマーク & MMLU-Pro 500 / MATH-500（Level 4--5）200 / SuperGPQA 300 \\
\hline
シード & $K = 3$（1, 2, 3．問題抽出・生成・同数決着に共通使用） \\
\hline
Self-Consistency & $k = 9$ \\
\hline
bootstrap反復 & $B = 10{,}000$ \\
\hline
多重比較補正 & Holm--Bonferroni（各実験のプール比較6件） \\
\hline
\end{tabular}
\end{table}

\begin{table}[htbp]
\caption{サービング・インフラの設定}
\label{tab:4.hp-infra}
\centering
\small
\begin{tabular}{|p{5.4cm}|p{7.6cm}|}
\hline
項目 & 値 \\
\hline
vLLM & \texttt{--enable-lora --max\_loras 8 --max\_lora\_rank 64 --max-model-len 32768}，TP=1 \\
\hline
動的ロード & \texttt{/v1/load\_lora\_adapter}（\texttt{VLLM\_ALLOW\_RUNTIME\_LORA\_UPDATING=True}） \\
\hline
学習ジョブ & Vertex AI Spot，n1-standard-8 + T4 $\times$1 \\
\hline
進化・評価ジョブ & Vertex AI Spot，a2-highgpu-1g（A100 40GB $\times$1），us-central1 \\
\hline
成果物保存 & GCS（Cloud Storage FUSEマウント），問題単位JSONL逐次追記による再開設計 \\
\hline
\end{tabular}
\end{table}

実験2（再学習，\S\ref{sec:4.6.3}）で表\ref{tab:4.hp-qlora}から変更した項目を表\ref{tab:4.hp-run002}に示す．学習関連の変更は忘却の抑制（rank・エポック・学習率の低減はLoRAの忘却がこれらとともに増加するという知見に基づく）とリプレイ長CoTの切り捨て防止（学習系列長）を，議論・集約関連の変更は\S\ref{sec:4.6.3}の処方をそれぞれ目的とする．

\begin{table}[htbp]
\caption{実験2の再学習で変更したハイパーパラメータ}
\label{tab:4.hp-run002}
\centering
\begin{tabular}{|l|l|l|}
\hline
項目 & 実験1（世代0） & 実験2（再学習） \\
\hline
学習データ & ペルソナ例60例 & ペルソナ例60例＋リプレイ36例（約38\%） \\
LoRA rank / $\alpha$ & 32 / 64 & 16 / 32 \\
エポック数 & 3 & 2 \\
学習率 & $2 \times 10^{-4}$ & $1 \times 10^{-4}$ \\
学習系列長上限 & 1{,}024トークン（フレームワーク既定） & 8{,}192トークン \\
議論の更新指示 & 標準（自由更新） & 条件付き（自己誤り特定時のみ変更可） \\
発言の提示 & 役割名付き & 匿名化＋提示順シャッフル \\
チーム集約 & 多数決 & logprob重み付き投票 \\
\hline
\end{tabular}
\end{table}

\section{実験結果}
\label{sec:4.6}

本節では実験を2段階に分けて報告する．\textbf{実験1}は\S\ref{sec:4.3}の比較条件による主実験であり，チーム化の効果と進化的最適化の寄与を測定するとともに，敗因の診断（損失分解）を行う．\textbf{実験2}は実験1の診断に基づく処方（能力保持リプレイ・検証型集約・議論プロトコル改良）を，実測ゲート方式で検証する追加実験である．

\subsection{実験1: 主実験の結果}
\label{sec:4.6.1}

表\ref{tab:4.res-main}に全条件の精度（$K=3$シード平均）を示す．いずれの条件も回答抽出失敗率は3\%以下であり，条件間の非対称性はない．

\begin{table}[htbp]
\caption{実験1の主要結果（精度，$K=3$シード平均）}
\label{tab:4.res-main}
\centering
\begin{tabular}{|l|c|c|c|}
\hline
条件 & MMLU-Pro & MATH-500 & SuperGPQA \\
\hline
1: ベースsolo & 0.639 & 0.830 & 0.436 \\
2: Self-Consistency@9 & 0.667 & 0.862 & \textbf{0.467} \\
3: 素の議論 & 0.598 & \textbf{0.870} & \textbf{0.467} \\
3$'$: プロンプトペルソナ議論 & 0.624 & 0.860 & 0.466 \\
4: gen-0 LoRAチーム & \textbf{0.684} & 0.770 & 0.407 \\
5: 進化後LoRAチーム & 0.683 & 0.813 & 0.399 \\
6: 進化後solo最良 & 0.661 & 0.698 & 0.364 \\
\hline
\end{tabular}
\end{table}

主要な検定結果（問題IDクラスタの符号反転置換検定，\S\ref{sec:4.4}）は以下のとおりである．

\begin{itemize}
\item \textbf{チーム化の効果（MMLU-Pro）}: LoRAペルソナチーム（条件4）はベース単体に対し$+4.3$pt $[+2.0, +6.6]$（$p=0.0004$），素の議論に対し$+8.4$pt相当の有意な優位を示した．
\item \textbf{生成回数を揃えた比較}: 進化後チーム（条件5）はSC@9に対しMMLU-Proで$+1.6$pt（非有意）だが，3ベンチマーク全体では$-1.8$pt $[-3.5, -0.2]$（$p=0.028$）と\textbf{有意に劣った}．
\item \textbf{進化の寄与}: 条件5と条件4の差はMATH-500の$+4.3$pt $[+1.9, +6.8]$（$p=0.0005$）のみ有意で，全体では$+0.2$pt（非有意）であった．
\item \textbf{ベンチマーク依存性}: 素の議論はMATH-500で$+4.2$pt（$p=0.0003$）・SuperGPQAで$+3.1$pt（$p=0.020$）と有意に有効だが，MMLU-Proでは$-4.1$pt（$p=0.001$）と有意に有害だった．
\end{itemize}

すなわち実験1の結論は「LoRAペルソナチームは知識推論（MMLU-Pro）でのみ有効であり，生成回数を揃えたSelf-Consistencyには総合で及ばない」である．

\subsection{実験1の診断: 損失分解}
\label{sec:4.6.2}

敗因を特定するため，各問題について「3体のうち少なくとも1体が正解した割合」（oracle上限）を計算し，チーム誤差を（i）個体能力の損失，（ii）集約の損失に分解した（表\ref{tab:4.oracle}）．

\begin{table}[htbp]
\caption{損失分解（seed 1）：oracle上限と集約損失}
\label{tab:4.oracle}
\centering
\begin{tabular}{|l|c|c|c|c|c|}
\hline
ベンチマーク & solo平均 & oracle & 多数決 & 議論後 & 集約損失 \\
\hline
SuperGPQA & 0.335 & \textbf{0.517} & 0.363 & 0.387 & $-0.130$ \\
MMLU-Pro & 0.673 & 0.800 & 0.702 & 0.714 & $-0.086$ \\
MATH-500 & 0.667 & 0.815 & 0.705 & 0.810 & $-0.005$ \\
\hline
\end{tabular}
\\[2pt]
{\footnotesize 全列seed 1の値．「多数決」は3体のsolo回答の多数決で，票が同数に割れた場合は固定シードの一様乱択でタイブレークする．}
\end{table}

3つの発見が得られた．第一に，\textbf{個体能力の毀損}である．ペルソナSFT後のsolo精度はベース単体に対しSuperGPQAで$-7.2$pt，MATH-500のLevel 5層で$-27$pt低下していた．診断用に収集した生成ログの分析から，ペルソナSFT（平均779字の短い応答のみで学習）がベースモデルの思考連鎖を中央値2{,}423字から1{,}188字へ圧縮しており，長い導出を要する問題で選択的に失敗することが確認された．第二に，\textbf{集約の損失}である．SuperGPQAでは毀損した3体ですらoracle上限が0.517とベース単体（seed 1で0.407）を$+11$pt上回るにもかかわらず，多数決＋議論はその38.7\%しか回収していない．正解が少数派に留まる問題では多数決は構造的に失敗するうえ，議論が正しい多数派を崩す事例も確認された（MMLU-Pro 500問中，議論前に多数派が正解だった問題のうち26問が議論後に不正解へ転じた）．一方MATH-500では集約損失がほぼゼロであり，検証可能な領域では正しい解が議論を通じて自然に採択されることを示す．第三に，\textbf{選抜ノイズ}である．進化ログの再分析により，適応度セット100問での測定は標準誤差$\pm 4.9$ptを持ち，同一個体の測定チーム精度が世代間で$\pm 4$〜$7$pt変動していた．毎世代の最大値選抜はノイズの上振れを拾う正のバイアス（勝者の呪い）を持ち，開発セットでの$+12$ptの改善が最終評価で$\pm 0$に消えた現象を説明する．さらに進化前後のアダプタを直接比較すると，$\Delta W$のFrobeniusノルム比は全役割で1.000であるだけでなく，\textbf{$\Delta W$同士のcosine類似度も全役割・全252モジュールで0.9996以上}（連結$\Delta W$で0.9999--1.0001）であった．すなわち進化後の最終チームのアダプタは，重みとしてはSFT直後のgen-0アダプタと実質的に同一である．演算子単位で測っても，突然変異1回でcos $=0.9997$，交叉1回でcos $=0.9998$，gen-1からgen-5までの4世代分でもcos $=0.9998$であり，本実験の進化演算子（ノルム保存の相対スケール変異と$\alpha \in [0.3, 0.7]$の内挿交叉）は重み空間をほとんど動かしていなかった．したがってMATH-500で観測された進化の寄与$+4.3$ptは，重み空間における意味のある移動（たとえば損傷成分の打ち消し）によるものとは考えにくく，固定された極小の重み摂動に対して生成過程が敏感に応答した結果——同一設定でも生成環境の違いにより$+6$pt級の系統差が生じるという\S\ref{sec:4.6.4}の観察と同種の感度——として説明するのが最も保守的である（第\ref{cha:5}章で詳述）．

\subsection{実験2: 処方の検証（ゲート方式）}
\label{sec:4.6.3}

実験1の診断に基づき，以下の処方を実測ゲート方式（各段階に事前定義した合格基準を設け，不合格の変更は採用しない）で検証した．

\begin{description}
\item[G0（能力保持リプレイの生成）] ベースモデル自身にMATH訓練split Level 4--5およびMMLU-Pro検証splitの問題を解かせ，正解が検証できた長い思考連鎖36例（平均6{,}503字）を採取した．
\item[G1（再学習と個体能力の検収）] 既存ペルソナ例60例にリプレイ36例（約38\%）を混合し，rank 16・2エポック・学習率$1 \times 10^{-4}$・系列長8{,}192トークンで3ペルソナを再学習した（run002）．solo精度はMMLU-Pro 0.665--0.695，MATH-500 0.775--0.845，SuperGPQA 0.405--0.440となり，全ベンチマークで事前基準（0.63/0.75/0.37）を満たした．\textbf{能力毀損は事前基準を満たす水準まで回復し，一部はベース単体を上回った}．
\item[G2（集約方式の選定）] 議論プロトコルに条件付き更新（自己の誤りを特定できた場合のみ回答変更を許可）と発言の匿名化を導入した上で，SuperGPQA 150問で集約方式を比較した．結果は多数決0.433・logprob重み付き投票0.440・GenSelect型裁定0.433であり，重み付き投票を採用した（裁定型の上積みは確認されなかった）．
\item[G3（チーム性能の確認）] 再学習チーム（以下\textbf{条件7$^\ddagger$}: リプレイ再学習$\times$条件付き議論$\times$匿名化$\times$重み付き投票）をseed 1の3ベンチマークで評価し，MMLU-Pro 0.754・MATH-500 0.855・SuperGPQA 0.430を得てゲートを通過した（この時点の比較対象は実験1の測定値であり，後述\S\ref{sec:4.6.4}の環境統制前である点に注意）．
\end{description}

\subsection{測定環境の系統差の発見と統制}
\label{sec:4.6.4}

実験2の初期評価（条件7$^\ddagger$を3シードで評価し実験1の測定値と比較）では，条件7$^\ddagger$がSC@9をプールで$+1.3$pt上回る（$p=0.053$；MMLU-Pro単独では$+5.5$pt相当）という境界的な結果が得られた．効果の維持を確認する目的でseed 4--6を追加したところ優位が消失し，さらにSC@9の絶対精度がシード群間で不連続に変化する（MMLU-Proで$0.66$--$0.69 \rightarrow 0.71$--$0.76$）という不自然なパターンが観察された．シード群の境界は評価用コンテナイメージの再ビルド時点と一致していた．

そこで同一問題・同一シード・同一生成設定でSC@9を新旧両イメージで実行する統制実験を行った結果，\textbf{同一の500問に対し旧イメージ0.688・新イメージ0.752と，イメージ差のみで$+6$pt級の系統差}が確認された（予測の一致率は86\%に留まり，生成自体が変化している）．vLLM本体のバージョンは固定（v0.8.5）であり，差はイメージ再ビルド時に更新された依存ライブラリに起因すると推定される．興味深いことに，この系統差は\textbf{ベースモデルの素の生成に強く現れ（$+8$pt級），LoRAアダプタを適用した生成にはほとんど現れない}（旧チーム＝条件5の実験1進化後チームのMMLU-Proは旧0.714/新0.712とほぼ不変．以下，条件5を「旧チーム」，条件7$^\ddagger$を「再学習チーム」と呼ぶ）．

この発見を受け，主要4条件（ベース単体・SC@9・旧チーム・再学習チーム）を\textbf{新イメージに統一して全面的に再測定}し，以降の比較はすべて同一環境の測定値のみで行う．なお実験1の内部比較（\S\ref{sec:4.6.1}--\ref{sec:4.6.2}）はすべて旧イメージ内の対応あり比較であるため結論は影響を受けないが，実験1と実験2の数値の直接比較は不能である．また同一seedを指定してもvLLMの生成は並列実行状態に依存して完全には再現しないことも確認しており，評価インフラの再現性は小型LLM評価の系統誤差要因として軽視できない．

\subsection{実験2: 最終評価（同一測定環境）}
\label{sec:4.6.5}

新イメージに統一した最終結果を表\ref{tab:4.res-v2}に示す．SC@9と条件7$^\ddagger$は6シード（1--6），ベース単体と旧チームは3シード（1--3）で評価した．

\begin{table}[htbp]
\caption{実験2の最終結果（同一測定環境，精度）}
\label{tab:4.res-v2}
\centering
\begin{tabular}{|l|c|c|c|c|}
\hline
条件 & MMLU-Pro & MATH-500 & SuperGPQA & シード数 \\
\hline
1: ベース単体 & 0.727 & 0.818 & 0.431 & 3 \\
2: SC@9 & \textbf{0.740} & \textbf{0.873} & \textbf{0.486} & 6 \\
5: 進化後チーム（実験1の構成） & 0.685 & 0.793 & 0.420 & 3 \\
7$^\ddagger$: 再学習チーム & 0.716 & 0.867 & 0.431 & 6 \\
\hline
\end{tabular}
\end{table}

主要な検定結果（問題IDクラスタの符号反転置換検定，$[\cdot,\cdot]$はクラスタブートストラップ95\%信頼区間，Holm補正後いずれも判定不変）:
\begin{itemize}
\item \textbf{主検定（対SC@9，6シード・固有問題4{,}683件）}: 条件7$^\ddagger$は$-3.49$pt $[-4.50, -2.56]$（$p<10^{-4}$）と\textbf{有意に劣った}．ベンチマーク別ではMMLU-Pro $-2.4$pt $[-3.6, -1.3]$（$p=0.0001$），MATH-500 $-0.8$pt $[-2.4, +0.7]$（$p=0.35$，差は検出されず），SuperGPQA $-5.5$pt $[-7.4, -3.7]$（$p<10^{-4}$）である\footnote{参考: （シード, 問題）プールのMcNemar正確検定では$-2.98$pt（勝ち223問，負け402問，$p<10^{-6}$）．点推定の差はクラスタ解析が各問題を等重み平均するのに対しプールが各観測を等重み平均することによる．}．
\item \textbf{処方の効果（対旧チーム，3シード・固有問題2{,}584件）}: 条件7$^\ddagger$は旧チームに対し$+2.68$pt $[+1.23, +4.13]$（$p=0.0003$，Holm補正後有意）と\textbf{一貫して優った}（MMLU-Pro $+3.7$pt・MATH-500 $+6.1$pt・SuperGPQA $\pm 0.0$pt）．診断に基づく処方（リプレイ再学習・条件付き議論・匿名化・重み付き投票）の効果は測定環境の統制後も頑健である．
\item \textbf{対ベース単体}: 条件7$^\ddagger$は$-0.32$pt $[-1.66, +1.04]$（$p=0.65$）で差は検出されず，事前定義した等価マージン$\pm 2$ptに対する同等性判定（90\%信頼区間$[-1.47, +0.81] \subset \pm 2$pt）により\textbf{同等性が成立}した．実験1で確認された「ペルソナSFTチームの能力毀損」（旧チームはベースに$-3.00$pt，$p=0.0001$で有意に劣る）はベース水準まで回復している．ベンチマーク別ではMATH-500のみ$+3.4$pt $[+0.8, +6.0]$（$p=0.010$）と有意に上回る．
\end{itemize}

なお「MATH-500でSC@9と同等」という表現については，差は検出されなかったものの90\%信頼区間$[-2.13, +0.45]$の下限が等価マージン$-2$ptをわずかに超えるため，\textbf{同等性は未確定}である（差がないとも$2$pt以上劣るとも断定できない）．

計算量の観点では，条件7$^\ddagger$は1問あたり6生成とSC@9の9生成より少ないが，議論の第2ラウンドは他エージェントの発話全文を入力に含むため入力トークン数は生成回数に比例しない．実トークン数の集計は\S\ref{sec:4.6.6}に示す．本比較は「生成回数がより少ない相手に負けた」ことを示すが，厳密なcompute-matched比較（総トークン・GPU時間の一致）ではない点に注意されたい．実験2の到達点は「処方によりチームを\textbf{ベース単体と同等（数学では有意に優位）}まで回復させたが，生成回数マッチのSelf-Consistencyには及ばない」である．

\subsection{追加分析: sharingの事後検証・カテゴリ別効果・計算量の実測}
\label{sec:4.6.6}

査読的観点からの再検証として，以下の3つの追加分析を行った．

\paragraph{(a) fitness sharingの事後検証 — 本実験の構成では選抜に作用しなかった}
\S\ref{sec:3.3}のfitness sharingは，役割内サブ集団が$K=2$の場合，行動距離の対称性（$d(a,b)=d(b,a)$）により2個体のsharing係数が常に同値となり，$\arg\max_c \phi_c \, s(c) = \arg\max_c \phi_c$，すなわち\textbf{役割内選抜に一切影響しない}．進化ログの全記録（6世代$\times$3役割 = 18ケース）を検証した結果，係数は全ケースで役割内同値であり，raw適応度（Shapley値）による順位とsharing適用後の順位は完全に一致していた．さらに18ケース中16ケースでは行動距離がniche半径$\sigma_{\mathrm{share}}=0.3$を超えており，係数自体が1.0（割引の不発動）であった．したがって\textbf{本実験で実際に働いた選択圧はShapley値のみ}であり，「fitness sharingにより多様性を維持した」という設計意図は$K=2$の構成では機能していない．この設計上の欠陥は本研究の重要な反省点であり，sharingを有効に作用させるには$K \geq 3$のサブ集団が必要である（Shapley値は負になり得るため，その場合は非負化などsharingとの整合にも別途注意を要する）．

\paragraph{(b) MMLU-Proのカテゴリ別分析 — 「検証可能性」仮説はベンチマーク内では再現しない}
実験1で観察された議論効果の反転（MATH-500・SuperGPQAで正，MMLU-Proで負）を「検証可能な領域では議論が有効」と解釈できるかを検証するため，MMLU-Proの14分野を定量系（math, physics, chemistry, engineering, computer science: 検算可能性が高い）と知識系（law, history, philosophy等）に事後的に分け，問題IDクラスタ解析で素の議論（条件3）のベース比の効果を分野別に測定した．結果は仮説の予測と\textbf{一致しなかった}: 定量系$-3.4$pt（$p=0.11$）・知識系$-4.7$pt（$p=0.002$）でグループ間差は非有意（$p=0.61$）であり，特にMMLU-Proのmathカテゴリでは議論により$-9.3$pt（$p=0.04$）と\textbf{悪化}した（MATH-500では$+4.2$ptの改善であるにもかかわらず）．すなわち，実験1で観察された反転は問題内容の「検証可能性」よりも\textbf{ベンチマーク単位の特性}（出題形式・難易度分布・選択肢の有無などの複合）に結びついており，本論文の「検証可能性による統一的説明」は\textbf{ベンチマーク間の観察に基づく仮説の域を出ない}．この分析はカテゴリ分けを事後に定義した探索的なものであるが，仮説の適用限界を示す証拠として報告する．

\paragraph{(c) 計算量の実測 — 「生成回数マッチ」の実態}
「生成回数を揃えた比較はcompute-matchedと言えるのか」という疑問に答えるため，LLM呼び出しの完全記録（入力メッセージ全文と生成全文）から，新環境の全呼び出しについてQwen3-4Bの実トークナイザで入出力トークン数を集計した（表\ref{tab:4.tokens}）．議論の第2ラウンドは他エージェントの発話全文を入力に含むため，条件7$^\ddagger$の入力はSC@9の約5倍に達する．一方，自己回帰生成の計算コストを支配する\textbf{出力トークンはSC@9が約2.8倍多く}，入出力合計でもSC@9（18{,}419トークン/問）が条件7$^\ddagger$（15{,}804トークン/問）を上回った．したがって本比較は厳密なcompute-matchedではないものの，\textbf{チーム側が計算量の点で優遇されていた比較ではなく}，SC@9は総トークンでもチームを上回る強いベースラインであった．本論文では両条件の関係を「生成回数を揃えた比較（チーム6生成 $<$ SC@9の9生成）」と表現し，トークン実測値を上記のとおり併記する．

\begin{table}[htbp]
\caption{計算量の実測（新環境・1問あたり平均トークン数）}
\label{tab:4.tokens}
\centering
\begin{tabular}{|l|c|c|c|c|}
\hline
条件 & 呼び出し数 & 入力トークン & 出力トークン & 合計 \\
\hline
7$^\ddagger$: 再学習チーム（議論2ラウンド） & 6.0 & 9{,}925 & 5{,}879 & 15{,}804 \\
2: SC@9 & 9.0 & 1{,}955 & 16{,}464 & 18{,}419 \\
\hline
\end{tabular}
\\[2pt]
{\footnotesize 条件7$^\ddagger$は6シード6{,}000問（36{,}000呼び出し）の平均．SC@9は条件2のみを含むジョブ（recheck・robust，計3{,}500問・31{,}500呼び出し，ベンチマーク構成比は本評価と同一）の平均．}
\end{table}
